\errorcontextlines9999
\documentclass[aspectratio=169,professionalfonts,fleqn]{beamer}

\usepackage{fontspec}
\usepackage{mathtools}
\usepackage{unicode-math}
\usepackage{csquotes}
\usepackage{accsupp}

\usepackage[lining]{FiraSans}

% do not like the lambda
\setmonofont{FiraCode}[
   Path = data/fonts/,
   Scale = 1,
   Extension = .ttf,
   Contextuals=Alternate,
   BoldFont={*-Bold},
   UprightFont={*-Regular},
]

\setmathfont{FiraMath-Regular}[
   Path = data/fonts/,
   Scale = 1.05,
   Extension = .otf
]

% TODO: find better/matching font? fira math is weird for greek stuff
\setmathfont[range={it/{greek,Greek},bfit/{greek,Greek}}]{lato}
\setoperatorfont\mathsf

\makeatletter

\def\todo#1{\textbf{\color{red}TODO:~\itshape#1}}
\def\Y{}
\def\I{}\def\S{}\def\K{}
\usepackage{xspace}

\def\InstallShortcuts{%
%
\def\textlambda{\lambda}%
\renewcommand\l[1][@main]{\ensuremath{{\color{##1}\symbf\textlambda}}}%
\def\L{\texorpdfstring{\ensuremath{\textlambda}}{\text{Lambda}}}%
\renewcommand\Y[1][@main]{\ensuremath{{\color{##1}\mathbf{Y}}}\xspace}%
\renewcommand\I[1][uulmaccent]{\ensuremath{{\color{##1}\mathbf{I}}}\xspace}%
\renewcommand\S[1][uulmaccent]{\ensuremath{{\color{##1}\mathbf{S}}}\xspace}%
\renewcommand\K[1][uulmaccent]{\ensuremath{{\color{##1}\mathbf{K}}}\xspace}%
\def\ifelse{\textsf{\textsl{ifelse}}\xspace}%
\def\pred{\textsf{\textsl{pred}}\xspace}%
\def\isZero{\textsf{\textsl{isZero}}\xspace}%
\def\T##1{\ifmmode\text{\texttt{##1}}\else\texttt{##1}\fi}%
\def\say##1{\enquote{##1}}%
\def\RED##1##2{\smash{\ensuremath{{}##1_{\rlap{\raisebox{3pt}{$\scriptstyle\color{uulmaccent}##2$}}}\mkern-.705mu{}}}}%
\def\IRED##1##2{\smash{\ensuremath{~~\;{}##1_{\raisebox{3pt}{$\scriptstyle\color{uulmaccent}##2$}}}~~}}%
\def\alphaconv{\RED=\alpha}%
\def\etaconv{\RED=\eta}%
\def\ialphaconv{\IRED=\alpha}%
\def\ibaseconv{\IRED={\phantom{\alpha}}}%
\def\ietaconv{\IRED=\eta}%
\def\betared##1{\RED={\beta_{\raisebox{-3pt}{\(\scriptscriptstyle##1\)}}}}%
\def\ibetared##1{\IRED={\beta_{\raisebox{-3pt}{\(\scriptscriptstyle##1\)}}}}%
\def\b[##1->##2]{\ensuremath{\T{{\color{uulmaccent}[##1\(\;\rightarrow\;\)##2]}}}}%
\def\@noxspace{}%
\xspaceaddexceptions{\@noxspace}%
\def\BR[##1->##2]{\ensuremath{\T{{\color{uulmaccent}[##1\@noxspace$\;\rightarrow\;$##2\@noxspace]}}}}%
% TODO: beta short red
}
\InstallShortcuts

\usepackage{xcolor}
% colors
\definecolor{uulmlogoblue}{HTML}{89a2b3}
\definecolor{uulmblue}{HTML}{7D9AAA}
\definecolor{uulmaccent}{HTML}{A9A28D}
\definecolor{green}{HTML}{56AA1C}
\definecolor{red}{HTML}{A32638}
\definecolor{orange}{HTML}{DF6D07}
\definecolor{blue}{HTML}{26547C}
\definecolor{accent}{HTML}{563C5C}

\colorlet{@main}{accent}

\preto\input@path{{lib/beamer-themes}{lib/tikzpingus/tex}{lib/fancyqr}{lib/algorithms}}
\usepackage{lib/fancyqr/fancyqr}
\usepackage{lib/algorithms/nice-algo}
\SetAlgoVlined
\RestyleAlgo{plain}

%---required packages & variable definitions------------------------------------
\usepackage{forest}
\usepackage{xcolor}
\usetikzlibrary{angles}
\definecolor{drawColor}{RGB}{64, 64, 64}
\newcommand{\circleSize}{0.165em}
\newcommand{\angleSize}{0.8em}
%-------------------------------------------------------------------------------
%---Define the style of the tree------------------------------------------------
\forestset{
	/tikz/mandatory/.style={
		circle,fill=drawColor,
		draw=drawColor,
		inner sep=\circleSize
	},
	/tikz/optional/.style={
		circle,
		fill=white,
		draw=drawColor,
		inner sep=\circleSize
	},
	featureDiagram/.style={
		for tree={
			text depth = 0,
			parent anchor = south,
			child anchor = north,
			draw = drawColor,
			rounded corners=2pt,
			edge = {draw=drawColor},
		}
	},
	/tikz/abstract/.style={
		fill = @main!15!white,
		draw = drawColor
	},
	/tikz/concrete/.style={
		fill = none,
		draw = drawColor
	},
	mandatory/.style={
		edge label={node [mandatory] {} }
	},
	optional/.style={
		edge label={node [optional] {} }
	},
	or/.style={
		tikz+={
			\path (.parent) coordinate (A) -- (!u.children) coordinate (B) -- (!ul.parent) coordinate (C) pic[fill=drawColor, angle radius=\angleSize]{angle};
		}
	},
	/tikz/or/.style={
	},
	alternative/.style={
		tikz+={
			\path (.parent) coordinate (A) -- (!u.children) coordinate (B) -- (!ul.parent) coordinate (C) pic[draw=drawColor, angle radius=\angleSize]{angle};
		}
	},
	/tikz/alternative/.style={
	},
	/tikz/placeholder/.style={
	},
	collapsed/.style={
		rounded corners,
		no edge,
		for tree={
			fill opacity=0,
			draw opacity=0,
			l = 0em,
		}
	},
	/tikz/hiddenNodes/.style={
		midway,
		rounded corners,
		draw=drawColor,
		fill=white,
		minimum size = 1.2em,
		minimum width = 0.8em,
		scale=0.9
	},
}
%-------------------------------------------------------------------------------

\def\pingulibpathprefix{lib/tikzpingus/tex/}
\usepackage[glows]{lib/tikzpingus/tex/tikzpingus}
\makeatletter
\usetikzlibrary{tikzmark,arrows.meta}
% The aux file is read inside the document group, where every \csname of a
% new mark creates a local \relax entry. With tens of thousands of
% remember-picture marks this exhausts TeX's save stack (save size), so we
% scope each entry; the actual definitions are global anyway.
\let\fp@savepicturepage\savepicturepage
\def\savepicturepage#1#2{\begingroup\fp@savepicturepage{#1}{#2}\endgroup}
\let\fp@pgfsyspdfmark\pgfsyspdfmark
\def\pgfsyspdfmark#1#2#3{\begingroup\fp@pgfsyspdfmark{#1}{#2}{#3}\endgroup}
\tikzset{
   @O/.style={overlay,remember picture},
   T/.style={uulmaccent!60!lightgray,font=\footnotesize\itshape\sffamily,line cap=round},
   lcr/.style={line cap=round}
}
\def\HStrut{\vphantom{\{tg}}% }<- wrong hl
\def\Snode#1{\tikzmarknode{#1}\HStrut}%
\def\TSnode#1{\subnode{#1}\HStrut}%
\def\bnode#1#2{\rnode{#1}{\hsGetStyle{basic}{#2}}}%
\def\rnode#1#2{\Snode{#1}#2\Snode{#1@}}
\def\TBM#1#2{\Snode{#1}#2\Snode{#1@}}
\def\hlcolor{@main}\def\hlopa{0.25}\def\hlhs{.4mm}% horizonal shift
\newcommand\hlrawbehindcodeunder[4][0pt]{\setbox0=\hbox{\pgfinterruptpicture\HStrut\endpgfinterruptpicture}\fill[\hlcolor,opacity=\hlopa,rounded corners=1pt] ([shift={(-\hlhs,-#1-1.6pt)}]#4#2) rectangle ([shift={(\hlhs,-#1+1.6pt)}]#4#2@); \coordinate (#3) at (#4#2@);
\path ([shift={(-\hlhs,-#1+.95\dp0)}]#4#2) -- ([shift={(\hlhs,-#1+.95\ht0)}]#4#2@) coordinate[pos=.5,yshift=1mm] (#3_up);
\path ([shift={(-\hlhs,-#1-.95\dp0)}]#4#2) -- ([shift={(\hlhs,-#1+.95\ht0)}]#4#2@) coordinate[pos=.5,yshift=-1mm-1.6pt] (#3_down); }
\newcommand\hlbehindcodeunder[3][0pt]{\hlrawbehindcodeunder[#1]{#2}{#3}{pic cs:}}
\newcounter{ULcounter}
\def\UL#1{\Snode{@ul@\theULcounter}#1\Snode{@ul@\theULcounter @}\tikz[@O]{\hlbehindcodeunder{@ul@\theULcounter}{@>}}\stepcounter{ULcounter}}

\def\G#1{\textcolor{gray}{#1}}
\def\Gopen{\textcolor{gray}{(}}
\def\Gclose{\textcolor{gray}{)}}
\newcommand\hlbehindcode[3][2pt]{\setbox0=\hbox{\pgfinterruptpicture\HStrut\endpgfinterruptpicture}\fill[\hlcolor,opacity=\hlopa,rounded corners=1pt] ([shift={(-\hlhs,-#1-.95\dp0)}]pic cs:#2) rectangle ([shift={(\hlhs,#1+.95\ht0)}]pic cs:#2@); \coordinate (#3) at ([yshift=.475\ht0]pic cs:#2@);
\path ([shift={(-\hlhs,#1+.95\dp0)}]pic cs:#2) -- ([shift={(\hlhs,#1+.95\ht0)}]pic cs:#2@) coordinate[pos=.5,yshift=1mm] (#3_up);\path ([shift={(-\hlhs,-#1-.95\dp0)}]pic cs:#2) -- ([shift={(\hlhs,-#1+.95\ht0)}]pic cs:#2@) coordinate[pos=.5,yshift=-1mm] (#3_down); }% hardcoded 1mm padding
% just underlines as there are no vert buffers:D

\newsavebox\pingubox


\let\hl\relax
\usepackage{lib/latex-underline/fancy-underline}

% Konfiguration der Unterstriche bei Links
\def\fancyul@depth{.965\p@}
\def\fancyul@contour{.2\p@}
\colorlet{tcbcolback}{white}
\def\fancyulbackground{tcbcolback}% so that it adapts with the beamer
% we must hook the background for the boxes

\colorlet{url@line@color}{uulmaccent}
\fancyulcolor{url@line@color!45!\fancyulbackground}
\newdimen\ul@surround \ul@surround=1.35\p@
% we make the auto detection of the depth bigger
\def\update@ul{\setbox\z@\hbox{{(j}}\edef\fancyul@thickness{\the\dimexpr.2\dp\z@}}
\def\rawlink#1{\leavevmode\kern-\ul@surround{\update@ul\kern\ul@surround\relax\fancyul{#1}}}
% define a link macro that can be used with href and fancyul
\DeclareRobustCommand*\link{\hyper@normalise\link@do}
\def\link@do#1#2{\leavevmode\kern-\ul@surround\href{#1}{\update@ul\kern\ul@surround\relax\fancyul{#2}}}% thinspace for some nicer padding
\def\xlink@color{hs@colors@lst@keywordA}
\DeclareRobustCommand*\Xlink{\hyper@normalise\Xlink@do}
\def\Xlink@do#1#2{\leavevmode\kern-\ul@surround\href{#1}{\update@ul\kern\ul@surround\relax\color{\xlink@color}\fancyul{#2}}}% thinspace for some nicer padding
\DeclareRobustCommand*\hlink{\hyper@normalise\hlink@do}
\def\hlink@do#1#2{\leavevmode\kern-\ul@surround\hyperlink{#1}{\update@ul\kern\ul@surround\relax\fancyul{#2}}}%
\usepackage{letltxmacro}
\AtBeginDocument{%
\LetLtxMacro\basicurl\url
% \def\url{\begingroup\catcode`{11\catcode`}11\obeyspaces\obeylines\url@do}%
% \def\url@do#1{\endgroup\link{#1}{#1}}%
\DeclareRobustCommand*\url{\hyper@normalise\url@do}
\def\url@do#1{\leavevmode\kern-\ul@surround\href{#1}{\update@ul\kern\ul@surround\relax\fancyul{#1}}}%
\long\def\beamer@ref#1{\hlink@do{#1}{{\beamer@origref*{#1}}}}% the extra set of braces is for soul so that it can work with braces
}
\def\abs#1{\left|#1\right|}
\def\N{\ensuremath{\mathbb{N}}}

\colorlet{@thegray}{uulmaccent!20!gray}
\def\info#1{\textcolor{@thegray}{\scriptsize\colorlet{@}{@thegray}#1}}
\def\HL#1#2{\T{\link{#1}{#2}}}
\def\DescHl{}
\def\Mix{15}\def\MixB{5}

\def\DescHeight{1.75cm}\def\DescText#1#2{\textbf{\vphantom{A}\smash{\DescHl\TLnk{#1}{#1}}}\nodepart{second}\parbox[c][\DescHeight]{.9\linewidth}{\footnotesize\DescHl\centering\T{#2}\endgraf}}
\def\RawDescText#1#2{\textbf{\vphantom{A}\smash{\DescHl#1}}\nodepart{second}\parbox[c][\DescHeight]{.9\linewidth}{\footnotesize\DescHl\centering\T{#2}\endgraf}}
\tikzset{@/.style={}}%
\def\BSep{5mm}
\def\RSep{9mm}

\def\EnumPhrase{\link{https://stackoverflow.com/questions/58761160/why-isnt-enum-typeclass-a-subclass-of-ord-typeclass}{Enum is no subclass of Ord?}}
\def\Yields{~~\textcolor{uulmaccent}{\faCaretRight}~~}%
% \gappto\OuterStuff{\onslide<#1|handout:#2>{#4}}
\def\ActiveOn#1#2#3{{\def\Mix{15}\def\MixB{5}\def\DescHl{}\only<3->{\def\Mix{5}\def\MixB{0}\def\DescHl{\color{gray}}}\only<#1|handout:#2>{\tikzset{K/.append style={black}}\def\Mix{15}\def\MixB{5}\def\DescHl{}}#3}}
\newcommand\Lnk[3][Prelude]{\T{\href{https://hackage.haskell.org/package/base-4.21.0.0/docs/#1.html\#v:#2}{#3}}}
\newcommand\LLnk[2][Prelude]{\T{\href{https://hackage.haskell.org/package/base-4.21.0.0/docs/#1.html\#v:#2}{\fancyul{#2}}}}
\newcommand\TLnk[3][Prelude]{\T{\href{https://hackage.haskell.org/package/base-4.21.0.0/docs/#1.html\#t:#2}{#3}}}
\newcommand\MLnk[2][Prelude]{\T{\href{https://hackage.haskell.org/package/base-4.21.0.0/docs/#1.html}{\fancyul{#2}}}}
\newcommand\TTLnk[2][Prelude]{\T{\href{https://hackage.haskell.org/package/base-4.21.0.0/docs/#1.html\#t:#2}{\fancyul{#2}}}}
\def\BLnk#1#2{\T{\href{https://hackage.haskell.org/package/base-4.21.0.0/docs/Control-Applicative.html\#v:#1}{#2}}}
\tikzset{K/.style={rectangle,draw,rounded corners=3.5pt,text width=3cm,minimum height=1.9cm,align=center,inner sep=5pt,rectangle split,rectangle split parts=2,rectangle split part fill={#1!\Mix!white,#1!\MixB!}}}

\def\CodeFileMarker#1{\CodeMarker{\llap{\faFileCodeO~\thinspace}#1}}
\def\CodeMarker#1{{\color{lightgray!75!gray}\sffamily\footnotesize#1}}
\colorlet{@filemarker}{lightgray!75!gray}
\let\curpath\@empty

\def\StartExcursion{%
   \begingroup
   \setbeamercolor{background canvas}{bg=gray!15}%
   \appto\PlaceFramenumberUpperLeft{\smash{\rlap{\tikz[remember picture,overlay]\node[below right,font=\tiny,gray!50,scale=.75] at(current page.north west) {excursion};}}}%
}
\def\EndExcursion{\endgroup}
\AtEndPreamble{\colorlet{tcbcolback}{btdl@color@background}}
\let\hl\relax

% manually load everything but sfmath
\usepackage{tikz,fontawesome,multicol,etoolbox,xstring,pbox}
\usetikzlibrary{calc,decorations.text,decorations.pathreplacing}
\usepackage[nolibs,nofonts,centerfoot]{lib/beamer-themes/dividing-lines/beamerthemedividing-lines}
\setbeamerfont{frametitle}{size=\Large,series=\mdseries}

\usepackage{lib/latex-lambda-calculus-visualizer/lc-visualizer}
\makeatletter
\tikzset{lc@abstraction@bracket/.append style={thin,opacity=.35},lc@node@replacement@collapsed/.append style={font=\ttfamily,execute at begin node={\vphantom{t}}},path image shift/.style={},
path image/.style 2 args={path picture={\node at ([path image shift]path picture bounding box.center) {\includegraphics[width=#2]{#1}};}}}
\let\fakemix\lc@colormix@doublevar

\lcSetup{duplicate-mixin={#1!83!yellow!60!lightgray!65!black}}
\newcommand*\reducearrow[3][]{\tikz[overlay,remember picture,uulmaccent!20!lightgray]{
   \draw[line cap=round,opacity=.5,line width=1mm] ([yshift=-.5mm]#2.north) to[edge node={node (@1) {}}] ([yshift=-.5mm]#2@.north);
   \path(#3.north) -- (#3@.north) coordinate[pos=.5] (@2);
   \draw[-Kite,bend right=10,#1] (@1) to (@2);
}}
\colorlet{lc@bracketcolor}{black}


\usepackage[super]{nth}

% patch it to include the sheet
\makeatletter
\def\sheet#1{\def\@sheet{#1}}
\let\@sheet\@empty% no \sheet: the title goes without a number
\let\oldtitle\title
\def\titlejoiner{.~}
\def\title#1{\oldtitle{\ifx\@sheet\@empty\else\@sheet\titlejoiner\fi#1}}
\colorlet{btdl@color@primary}{@main}%
% \definecolor{btdl@color@secondary}{RGB}{#2}%
% \definecolor{btdl@color@alerted}{RGB}{#3}%

\newcounter{BLcounter}
% \newcommand<>\BL[2]{\Snode{@bl@\theBLcounter}#2\Snode{@bl@\theBLcounter @}{\def\hlcolor{#1}\def\hlopa{.15}\only#3{\tikz[@O]{\hlbehindcode{@bl@\theBLcounter}{@>}}\stepcounter{BLcounter}}}}
\newcommand<>\BL[2]{#2}%

\usepackage{polyglossia}
\setdefaultlanguage[variant=american]{english}

\InstallShortcuts
\def\optional{\textcolor{uulmaccent}{\textit{(optional)}}\xspace}

% \titlehead{FP WT22/23}
\subtitle{Functional Programming, Winter Term 2026/27}
\institute{Software Engineering and Programming Languages}

\author[\href{mailto:florian.sihler@uni-ulm.de}{Florian Sihler}]{\href{mailto:florian.sihler@uni-ulm.de}{Florian Sihler}}
\outro{Ulm, \insertdate}
\newenvironment{tasks}{\begingroup\setbeamertemplate{enumerate item}{\itshape\alph{enumi})}\settowidth\leftmargini{xx)}%
\addtolength\leftmargini{\labelsep}\begin{enumerate}}{\end{enumerate}\endgroup}
\DeclareSymbolFont{cmsymbols}{OMS}{cmsy}{m}{n}
\SetSymbolFont{cmsymbols}{bold}{OMS}{cmsy}{b}{n}
\DeclareSymbolFontAlphabet{\mathcal}{cmsymbols}
\AtBeginDocument{\let\setminus\relax
\DeclareMathOperator\setminus{\textbackslash\,}}

\def\F{\ensuremath{\mathcal{F}}}
\newsavebox\pinguexplainbox

% shared by sheet-global and session-global.tex: magic-haskell, lc-visualizer and the old haskell-listings interface
\chardef\magiccompatatcode=\catcode`\@
\makeatletter
\preto\input@path{{lib/magic-haskell/}{lib/magic-haskell/xlistings/}{lib/magic-haskell/xlistings/langs/}{lib/latex-lambda-calculus-visualizer/}}
\RequirePackage[language=mhs]{magic-haskell}
\@ifpackageloaded{lc-visualizer}{}{\RequirePackage{lc-visualizer}}
\makeatother
\RequirePackage{xstring}

% compatibility with the old haskell-listings commands used on sheets and slides
\makeatletter
% the default colors of haskell-listings (the documents set keywordA-C themselves)
\lstcolorlet{literals}{uulmaccent!90!brown}
\lstcolorlet{numbers}{blue!70!black!80!teal}
\lstcolorlet{comments}{lightgray!75!gray}
\def\T{\texttt}
\lstset{autogobble}
% the old string delimiters (for strings spanning several inline snippets)
\lstdefinelanguage{lSheetMHS}{
   language=lMagicHaskell,
   moredelim=[is][{\color{xlst@colors@lst@literals}}]{:':>}{<:':},
   % from haskell-listings: grayed-out context lines and the emoji placeholder
   moredelim=[il][{\color{lightgray}}]{go>},
   add to literate={:Penguin:}{{{\emoji{penguin}}}}1
}
\xlstmintedmap{mhs}{lSheetMHS}
\xlstmintedmap{haskell}{lSheetMHS}% keeps the ligatures haskell-listings had
\newcommand\ihs[1][]{\bmhs[language=lSheetMHS,#1]}
\def\hsblacklistlinenumbers{\xlstblacklistlinenumbers}
\lstdefinelanguage{lSheetGHCi}{
   language=lGHCi,
   moredelim=[is][{}]{out->}{<-out},
   moredelim=**[il][{{\Large\strut}}]{in>},
   add to literate={ghci-Text.Regex.Posix>}{{{\mhs@nocopy{\mhsPrompt{ghci Text.Regex.Posix>}}}}}{22}
}
\xlstmintedmap{ghci}{lSheetGHCi}
\lstdefinelanguage{lText}{
   add to literate={:CRLF:}{{{\scriptsize\textcolor{lightgray}{\textbackslash r\textbackslash n}}}}3
}
\xlstmintedmap{text}{lText}
\newcommand\itxt[1][]{\bmhs[language=lText,#1]}
\colorlet{@rawtypecolor}{lightgray}
\def\@typeShift#1{\raisebox{.75pt}{\T{\hbox{#1}}}}
\robustify\@typeShift
\def\TypeSubstitute#1{\StrSubstitute{#1}{||>}{{\@typeShift{|\kern-3.5pt|>}}}[\@tmp]\StrSubstitute{\@tmp}{|>}{{\@typeShift{|>}}}[\@tmp]\StrSubstitute{\@tmp}{->}{{\ensuremath{\rightarrow}}}[\@tmp]\StrSubstitute{\@tmp}{=>}{{\ensuremath{\Rightarrow}}}[\@tmp]}
\robustify\TypeSubstitute
\def\RawTypePrint#1{\textcolor{@rawtypecolor}{\TypeSubstitute{#1}\BeginAccSupp{ActualText={}}\footnotesize\textsf{\@tmp}\EndAccSupp{}}}
% minted looks as with haskell-listings: plain (no frame, no line numbers), fixed columns,
% '-- !*...*!' escapes; text blocks never get line numbers
\xlstsetmintedstyle{plain}
\lst@AddToHook{InitVars}{\lst@width.6em\relax}% haskell-listings used basewidth=.6em (magic-haskell: the glyph width)
\apptocmd\xlst@minted@setup{\lstset{columns=fixed,lineskip=.33pt,escapeinside={--\ !*}{*!},keepspaces=false}\ifstrequal{#3}{text}{\lstset{numbers=none}}{}}{}{\PackageWarning{magic-compat}{Could not patch minted}}
\colorlet{prompt@color}{gray}
\def\HaskellPrompt#1{\textcolor{prompt@color}{#1}}
\colorlet{haskell@escapecolor}{uulmaccent!70!black}
\def\HaskellListingsEscape#1{\textcolor{haskell@escapecolor}{#1}}
\def\mhsEscape{\HaskellListingsEscape}% slides restyle \HaskellListingsEscape
\def\hsComment#1{{\xlstGetStyle{comments}#1}}
\def\hsKw#1{{\xlstGetStyle{keywordA}#1}}
\def\hsNum#1{{\xlstGetStyle{numbers}#1}}
\def\MtlBase{https://hackage.haskell.org/package/mtl/docs}
\def\ParsecBase{https://hackage.haskell.org/package/parsec/docs}
\def\ArrayBase{https://hackage.haskell.org/package/array/docs}
\def\ByteStringBase{https://hackage.haskell.org/package/bytestring/docs}
\def\QuickCheckBase{https://hackage.haskell.org/package/QuickCheck/docs}
\def\HspecBase{https://hackage.haskell.org/package/hspec/docs}
\def\HspecExpBase{https://hackage.haskell.org/package/hspec-expectations/docs}
\def\mhsPrompt{\HaskellPrompt}% slides locally redefine \HaskellPrompt to hide prompts
\def\HsBase{https://hackage.haskell.org/package/base/docs}
\def\blacklst{\appto\xlst@styles@lst@basic{\color{black}}}
\AtBeginDocument{\colorlet{hs@colors@lst@keywordA}{xlst@colors@lst@keywordA}}% used by \Xlink in shared.tex
\makeatother
\catcode`\@=\magiccompatatcode % restore the catcode of the caller

% \NoteRef{key}: speaker notes that live in notes/notes.tex (private), shown only by the note builds of the makefile
% Speaker notes kept outside the slides: notes/notes.tex holds one DefineNote entry per note (key and text), the frames
% hold a NoteRef with the key. Without notes/notes.tex (the public copy) every \NoteRef expands to nothing.
% tools/extract-notes moves new notes out of the slides.
% This file leaves the catcode of @ as it finds it (session-global.tex inputs it while @ is a letter).
\expandafter\chardef\csname LectureAtCatcode\endcsname=\catcode`\@
\makeatletter
\newcommand\DefineNote[2]{\long\expandafter\gdef\csname lecture@note@#1\endcsname{#2}}
% \NoteRef{key} is \note{text}, if the key is defined
\newcommand\NoteRef[1]{\ifcsname lecture@note@#1\endcsname
   \expandafter\expandafter\expandafter\note\expandafter\expandafter\expandafter{\csname lecture@note@#1\endcsname}%
\else\LectureNoteMissing{#1}\fi}
% a key without text only warns, and only if notes/notes.tex is there (the public copy has none)
\newif\iflecture@notesfile
\newcommand\LectureNoteMissing[1]{\iflecture@notesfile\PackageWarningNoLine{lecture}{No text for the note #1 in notes/notes.tex}\fi}
% the notes are read when the document starts, so that their text is read with the same catcodes as in the frames
\AtBeginDocument{\IfFileExists{notes/notes.tex}{\lecture@notesfiletrue\input{notes/notes.tex}}{}}

% The note builds (\LectureNotes defined, see the makefile) put the notes beside every slide and also into pdfpc
% annotations. All other builds leave the notes out, as beamer does by default: they show no notes at all.
\ifdefined\LectureNotes
% the pdfpc package was written for pdfLaTeX, LuaTeX names the primitive differently
\ifdefined\pdfannot\else\protected\def\pdfannot{\pdfextension annot }\fi
\usepackage{letltxmacro}
\usepackage{beamer-latex-pdfpc-notes}
% the plain text of a note for pdfpc: its markers become lines and labels
\makeatletter
\LetLtxMacro\LecturePdfpcPlain\note
% the annotation names its author, and the author of this deck is a link
\newcommand\LecturePdfpcAnnot[1]{\begingroup\def\@author{Speaker notes}\LecturePdfpcPlain{#1}\endgroup}
\ExplSyntaxOn
\tl_new:N \l__lecture_pdfpc_tl
\NewDocumentCommand\LectureNoteForPdfpc{m}
   {
      \tl_set:Nn \l__lecture_pdfpc_tl {~ #1}
      % commands that would not survive the expansion into the annotation become their plain text;
      % a parenthesis is escaped, so that an unbalanced one cannot end the string of the annotation
      \regex_replace_all:nnN { \c{ihs} \cB. (.*?) \cE. } { \1 } \l__lecture_pdfpc_tl
      \regex_replace_all:nnN { \c{textbackslash} \s* } { \cO\\ \cO\\ } \l__lecture_pdfpc_tl
      \regex_replace_all:nnN { \c{\&} } { \cO\& } \l__lecture_pdfpc_tl
      \regex_replace_all:nnN { \c{\#} } { \cO\# } \l__lecture_pdfpc_tl
      \regex_replace_all:nnN { ([\(\)]) } { \cO\\ \1 } \l__lecture_pdfpc_tl
      % pdfpc renders markdown: the main message in bold, then the points as a list, then one paragraph per block
      % the question that leads in, in italics above the main message
      \regex_replace_once:nnN { \s TRANSITION : \s (.*?) (\s (?: MAIN | POINT | SUB | NEXT | TIP | ASK | SHOW | EXTRA | DETOUR | ANECDOTE | GAG | READ | WHY\ IT\ MATTERS ) : | \Z ) }
         { _\1_ \c{\\} \2 } \l__lecture_pdfpc_tl
      \regex_replace_once:nnN { \s MAIN : \s (.*?) (\s (?: POINT | SUB | NEXT | TIP | ASK | SHOW | EXTRA | DETOUR | ANECDOTE | GAG | READ | WHY\ IT\ MATTERS ) : | \Z ) }
         { **\1** \c{\\} \2 } \l__lecture_pdfpc_tl
      \regex_replace_all:nnN { \s POINT : \s } { \c{\\}-\  } \l__lecture_pdfpc_tl
      \regex_replace_all:nnN { \s SUB : \s } { \c{\\}\ \ \ \ -\  } \l__lecture_pdfpc_tl
      \regex_replace_all:nnN { \s ASK : \s } { \c{\\}\c{\\}**Ask:**\  } \l__lecture_pdfpc_tl
      \regex_replace_all:nnN { \s SHOW : \s } { \c{\\}\c{\\}**Show:**\  } \l__lecture_pdfpc_tl
      \regex_replace_all:nnN { \s TIP : \s } { \c{\\}\c{\\}Tip:\  } \l__lecture_pdfpc_tl
      \regex_replace_all:nnN { \s EXTRA : \s } { \c{\\}\c{\\}Extra:\  } \l__lecture_pdfpc_tl
      \regex_replace_all:nnN { \s WHY\ IT\ MATTERS : \s } { \c{\\}\c{\\}Why\ it\ matters:\  } \l__lecture_pdfpc_tl
      \regex_replace_all:nnN { \s DETOUR : \s } { \c{\\}\c{\\}**Detour:**\  } \l__lecture_pdfpc_tl
      \regex_replace_all:nnN { \s ANECDOTE : \s } { \c{\\}\c{\\}Anecdote:\  } \l__lecture_pdfpc_tl
      \regex_replace_all:nnN { \s GAG : \s } { \c{\\}\c{\\}Gag:\  } \l__lecture_pdfpc_tl
      \regex_replace_all:nnN { \s NEXT : \s } { \c{\\}\c{\\}Next:\  } \l__lecture_pdfpc_tl
      \regex_replace_all:nnN { \s READ : \s } { \c{\\}\c{\\}Read:\  } \l__lecture_pdfpc_tl
      \regex_replace_all:nnN { \s ANSWER : \s } { \c{\\}\c{\\}Answer:\  } \l__lecture_pdfpc_tl
      \regex_replace_all:nnN { \s SOURCE : \s } { \c{\\}Source:\  } \l__lecture_pdfpc_tl
      \exp_args:NV \LecturePdfpcAnnot \l__lecture_pdfpc_tl
   }
\ExplSyntaxOff
% the note pages themselves, beside every slide
\input{setup/notes}
\fi
\catcode`\@=\LectureAtCatcode


\newcommand<>\CodeFileMarkerAttach[2][]{\only#3{\CodeFileMarker{\attachfilesetup{color=@filemarker}\fancyulcolor{@filemarker}\textattachfile{\curpath #1#2}{\rawlink{#2}}}}}
\newcommand<>\CodeFileMarkerAttachRaw[2][]{\only#3{\attachfilesetup{color=@filemarker}\StrSubstitute{#2}{_}{\_}[\name]\textattachfile{\curpath #1#2}{\name}}}
\newcommand<>\CodeFileAttach[2][]{\only#3{\attachfilesetup{color=@filemarker}\StrSubstitute{#2}{_}{\_}[\name]\textattachfile{\curpath #1#2}{\expandafter\rawlink\expandafter{\name}}}}
\newcommand<>\ShowInTheWeb[2][.25mm]{\ShowInTheWebWith#3[#1]{#2}{run in browser}}
\newcommand<>\ShowInTheWebWith[3][.25mm]{\only#4{\smash{\raisebox{#1}{\color{lightgray}\sffamily\tiny\fancyulcolor{lightgray}\link{#2}{{\raisebox{1pt}{\tikz[scale=.5]{\draw[rounded corners=.33pt] (0,0) -- (0,.25) -- (.3,.125) -- cycle;}}} #3}}}}\vspace*{-#1}}

\usepackage{attachfile2}
\attachfilesetup{%
    author={Florian Sihler},%
    color=0.2 0.3647 0.4941, icon=Tag, mimetype=text/plain%
}
\colorlet{pingu-board-back}{gray!85!red!10!white}
\newcommand\MakeThePinguExplainIt[3][]{%
\setbox\pinguexplainbox=\hbox{\tikzpicture\pinguexplain{#1}{#2}{#3}\endtikzpicture}}
\def\pinguexplain#1#2#3{
  \pingu[body type=legacy,body=pingu@black!70!pingu@main,glasses,left wing wave,staff left,right wing grab,cup=pingu@green!80!pingu@black,left item angle=-10,left eye wink,staff left length=23mm,cap=pingu@bronze!50!pingu@black!70!pingu@main, cap extra height=1.125mm,name=pingu,#2]

  \node[above right=-2.66mm,fill=pingu-board-back,draw=pingu@bronze!50!pingu@black!70!pingu@main!80!black,line width=.6mm,rounded corners=2pt,text width=5.25cm,minimum height=1.5cm,text=black,inner sep=8pt,align=justify,#1] at (pingu-staff-left-top) {\topsep=0pt\ignorespaces#3};}% TODO: LstBoardStyle

\newsavebox\tempboxA
\usepackage{forest}

% maybe?
% \appto\xlst@styles@lst@basic{\color{@main!15!black}}%
\lstcolorlet{keywordA}{@main!70!green}
\lstcolorlet{keywordB}{@main!40!green}
\lstcolorlet{keywordC}{@main!85!green}
\lstcolorlet{numbers}{purple!60!@main!90!red}
% \def\xlst@styles@lst@keywordA{\@declaredcolor{xlst@colors@lst@keywordA}\sbfamily}


\def\tokenchain#1#2#3{%
  \foreach[count=\i, remember=\i as \li (initially 0)] \elem in {#3}{%
    \ifnum\li>0
      \node[inner xsep=0pt,right] (#1\i) at (#1\li.east) {\!\strut\expandafter\ihs\expandafter{\elem}\!};
    \else
      \node[inner xsep=0pt,right] (#1\i) at (#2) {\!\strut\expandafter\ihs\expandafter{\elem}\!};
    \fi
    \xdef\@lasti{\i}%
  }%
}

% #1 name, #2 startpos, #3 chain, #4 type
\def\tokenchaintype#1#2#3#4{%
  \tokenchain{#1}{#2}{#3}%
  \node[inner xsep=0pt,right] (#1type) at (#1\@lasti.east) {\strut\small\type{#4}};
}
\preto\input@path{{lib/code-animation}}
\usepackage{lib/code-animation/code-animation}
\def\ca@Strut{\vphantom{\xlst@font@fs@normal\xlst@styles@lst@basic\strut}}

% no space here!
\xspaceaddexceptions{\Gclose}%
\def\lcnegskip{\vspace*{-.85\bigskipamount}}
